% GENERATED FILE. Edit canonical lessons, then run npm run generate:books.
% canonical-source-hash: fnv1a64:f0184b0a391f17ea
% canonical-lessons: MW-C123-takiyo, MW-C123-chadar, MW-C123-dari, MW-C123-pardo, MW-C123-arsi

\chapter{Pillow, Sheet, Cotton rug, Curtain, Mirror}
\label{ch:mw-pillow-sheet-cotton-rug-curtain-mirror}
\hlchaptermodality{123}

\begin{chapteropening}
\textbf{By the end of this chapter:} \emph{I can say a pillow, a sheet, a cotton rug, a curtain and a mirror.}

Complete the last lesson of chapter 123: I can say a pillow, a sheet, a cotton rug, a curtain and a mirror.
\end{chapteropening}

\section[\texorpdfstring{\mw{तकियो} (takiyo)}{तकियो (takiyo)}]{\mw{तकियो} (takiyo) — a pillow}

\label{lesson:MW-C123-takiyo}



\begin{quote}
Before the new one: say the Marwadi for a quilt, then the Marwadi for a mattress.
\end{quote}

\subsection*{You'll want to know: \mw{तकियो}}
\textbf{\mw{तकियो}} — \emph{takiyo} — \textquotedblleft{}a pillow\textquotedblright{}.

\begin{cousinweb}
\textbf{\mw{तकियो}} (\emph{takiyo}) is a loan, usually traced to Persian \emph{takya}, \textquotedblleft{}pillow, support\textquotedblright{}. Persian took it from Arabic, from a root meaning \textquotedblleft{}to lean\textquotedblright{} (Lane, root \emph{w-k-ʾ}).
\end{cousinweb}

\begin{grammarlens}[title={What you've built}]
One more word for this chapter.
\end{grammarlens}

\subsection*{Your turn}
\begin{itemize}
\raggedright
  \item \emph{Say it:} \emph{takiyo}
  \item \emph{Say it:} \emph{takiyo}, once more
  \item \emph{Say it:} say \emph{takiyo}
  \item \emph{Recall:} say the Marwadi for a quilt, then the Marwadi for a mattress, then say \emph{takiyo} again
\end{itemize}

\begin{tcolorbox}[breakable,colback=teal!4,colframe=teal!35!black,title={Before you move on}]
What does \mw{तकियो} mean? (\textquotedblleft{}A pillow\textquotedblright{}.) Say it once more.
\end{tcolorbox}
\section[\texorpdfstring{\mw{चादर} (chādar)}{चादर (chādar)}]{\mw{चादर} (chādar) — a sheet}

\label{lesson:MW-C123-chadar}



\begin{quote}
Before the new one: say the Marwadi for a mattress, then the Marwadi for a pillow.
\end{quote}

\subsection*{You'll want to know: \mw{चादर}}
\textbf{\mw{चादर}} — \emph{chādar} — \textquotedblleft{}a sheet\textquotedblright{}.

\begin{cousinweb}
\textbf{\mw{चादर}} (\emph{chādar}) is a loan, usually traced to Persian \emph{čādar}, \textquotedblleft{}sheet, veil\textquotedblright{}.
\end{cousinweb}

\begin{grammarlens}[title={What you've built}]
One more word for this chapter.
\end{grammarlens}

\subsection*{Your turn}
\begin{itemize}
\raggedright
  \item \emph{Say it:} \emph{chādar}
  \item \emph{Say it:} \emph{chādar}, once more
  \item \emph{Say it:} say \emph{chādar}
  \item \emph{Recall:} say the Marwadi for a mattress, then the Marwadi for a pillow, then say \emph{chādar} again
\end{itemize}

\begin{tcolorbox}[breakable,colback=teal!4,colframe=teal!35!black,title={Before you move on}]
What does \mw{चादर} mean? (\textquotedblleft{}A sheet\textquotedblright{}.) Say it once more.
\end{tcolorbox}
\section[\texorpdfstring{\mw{दरी} (darī)}{दरी (darī)}]{\mw{दरी} (darī) — a cotton rug}

\label{lesson:MW-C123-dari}



\begin{quote}
Before the new one: say the Marwadi for a pillow, then the Marwadi for a sheet.
\end{quote}

\subsection*{You'll want to know: \mw{दरी}}
\textbf{\mw{दरी}} — \emph{darī} — \textquotedblleft{}a cotton rug\textquotedblright{}.

\begin{cousinweb}
\textbf{\mw{दरी}} (\emph{darī}) is a tadbhava, a word inherited and reshaped by sound change. It comes from \emph{*darikā}, \textquotedblleft{}split cane, mat\textquotedblright{}, a form CDIAL reconstructs from its descendants. CDIAL gives no Marwari form in this entry (CDIAL 6194); the match is through Hindi \emph{darī}. It is built on the root \emph{dṛ}, \textquotedblleft{}to split\textquotedblright{}.
\end{cousinweb}

\begin{grammarlens}[title={What you've built}]
One more word for this chapter.
\end{grammarlens}

\subsection*{Your turn}
\begin{itemize}
\raggedright
  \item \emph{Say it:} \emph{darī}
  \item \emph{Say it:} \emph{darī}, once more
  \item \emph{Say it:} say \emph{darī}
  \item \emph{Recall:} say the Marwadi for a pillow, then the Marwadi for a sheet, then say \emph{darī} again
\end{itemize}

\begin{tcolorbox}[breakable,colback=teal!4,colframe=teal!35!black,title={Before you move on}]
What does \mw{दरी} mean? (\textquotedblleft{}A cotton rug\textquotedblright{}.) Say it once more.
\end{tcolorbox}
\section[\texorpdfstring{\mw{पड़दो} (paṛdo)}{पड़दो (paṛdo)}]{\mw{पड़दो} (paṛdo) — a curtain}

\label{lesson:MW-C123-pardo}



\begin{quote}
Before the new one: say the Marwadi for a sheet, then the Marwadi for a cotton rug.
\end{quote}

\subsection*{You'll want to know: \mw{पड़दो}}
\textbf{\mw{पड़दो}} — \emph{paṛdo} — \textquotedblleft{}a curtain\textquotedblright{}.

\begin{cousinweb}
\textbf{\mw{पड़दो}} (\emph{paṛdo}) is a loan, usually traced to Persian \emph{parda}, \textquotedblleft{}curtain, veil\textquotedblright{}. Marwadi says the \emph{r} as \emph{ṛ}.
\end{cousinweb}

\begin{grammarlens}[title={What you've built}]
One more word for this chapter.
\end{grammarlens}

\subsection*{Your turn}
\begin{itemize}
\raggedright
  \item \emph{Say it:} \emph{paṛdo}
  \item \emph{Say it:} \emph{paṛdo}, once more
  \item \emph{Say it:} say \emph{paṛdo}
  \item \emph{Recall:} say the Marwadi for a sheet, then the Marwadi for a cotton rug, then say \emph{paṛdo} again
\end{itemize}

\begin{tcolorbox}[breakable,colback=teal!4,colframe=teal!35!black,title={Before you move on}]
What does \mw{पड़दो} mean? (\textquotedblleft{}A curtain\textquotedblright{}.) Say it once more.
\end{tcolorbox}
\section[\texorpdfstring{\mw{आरसी} (ārsī)}{आरसी (ārsī)}]{\mw{आरसी} (ārsī) — a mirror}

\label{lesson:MW-C123-arsi}



\begin{quote}
Before the new one: say the Marwadi for a cotton rug, then the Marwadi for a curtain.
\end{quote}

\subsection*{You'll want to know: \mw{आरसी}}
\textbf{\mw{आरसी}} — \emph{ārsī} — \textquotedblleft{}a mirror\textquotedblright{}.

\begin{cousinweb}
\textbf{\mw{आरसी}} (\emph{ārsī}) is a tadbhava, a word inherited from Sanskrit and reshaped by sound change. It comes from \emph{ādarśa}, \textquotedblleft{}mirror\textquotedblright{}. CDIAL gives no Marwari form in this entry (CDIAL 1143); the match is through Hindi \emph{ārsī} and Gujarati \emph{ārsī}. It is built on the root \emph{dṛś}, \textquotedblleft{}to see\textquotedblright{}.
\end{cousinweb}

\begin{grammarlens}[title={What you've built}]
One more word for this chapter.
\end{grammarlens}

\subsection*{Your turn}
\begin{itemize}
\raggedright
  \item \emph{Say it:} \emph{ārsī}
  \item \emph{Say it:} \emph{ārsī}, once more
  \item \emph{Say it:} say \emph{ārsī}
  \item \emph{Recall:} say the Marwadi for a cotton rug, then the Marwadi for a curtain, then say \emph{ārsī} again
\end{itemize}

\begin{tcolorbox}[breakable,colback=teal!4,colframe=teal!35!black,title={Before you move on}]
What does \mw{आरसी} mean? (\textquotedblleft{}A mirror\textquotedblright{}.) Say it once more.
\end{tcolorbox}
